% =====================================================================
%  3 -- WHERE Y_e COMES FROM: THE WEAK SECTOR, QUANTITATIVELY.
%
%  Source: tier3.md 0.7 intro (2213-2216), 0.7.1 (2218-2263),
%          0.7.2 (2265-2321), 0.7.3 (2323-2350);
%          plus self-check Q9 from 0.10 (2736-2738) -> sec:selfcheck-weak.
%  Headings covered: "0.7 Where Ye comes from -- the weak sector,
%          quantitatively" (-> \section),
%          "0.7.1 The identity, and what can move it",
%          "0.7.2 Electron capture in a degenerate plasma -- why it is
%                 (nearly) a one-way ratchet",
%          "0.7.3 The measured Ye carriers".
%  Re-homings / rewordings: this section is MOVED UP from source 0.7 to
%          report section 3, i.e. it now precedes the collapse section
%          (source 0.4-0.5) and the observable end (source 0.6).
%          Consequently the opening sentence "SS0.3-0.6 established that
%          Ye is the quantity" is reworded to a mixed past/forward form
%          (TeX comment at the spot).  "the runaway of SS0.5.2" is left
%          as a forward \ref (fine).  Blockquote 2228-2229 -> invariant.
% =====================================================================

\section{Where \texorpdfstring{\Ye}{Ye} comes from --- the weak sector, quantitatively}
\label{part:weak}

% [edition] reworded cross-ref: "SS0.3-0.6 established that Ye is the quantity."
\S\ref{part:siburn} established, and \S\ref{part:collapse}--\S\ref{part:observables}
will establish, that \Ye{} is the quantity. This section is where it actually
changes, which is the sector this project's invariants protect most fiercely.

% =====================================================================
\subsection{The identity, and what can move it}
\label{sec:yeidentity}

From the definitions (Tier~0 \S II.5),
%
\begin{equation}
\Ye \;=\; \sum_i Z_i Y_i \;=\; \sum_i \frac{Z_i}{A_i}X_i ,
\label{eq:yeidentity}
\end{equation}
%
and the key structural fact:

\begin{invariant}
\textbf{No strong or electromagnetic reaction can change \Ye.} They conserve
$Z$ and $A$ separately, so $\sum Z_i\,\dd Y_i = 0$ identically for every
strong/EM column of $\nu$.
\end{invariant}

Nullity $\ge 2$ follows by construction (strong and electromagnetic reactions
conserve $Z$ and $A$ separately; \citealt[\S I]{langanke2003a}); that the left
null space of $\nu$ restricted to strong/EM columns is \emph{exactly}
2-dimensional and equals $\mathrm{span}\{A, Z\}$ --- the strong subgraph is
connected --- is measured in Tier~2 \S I.2.1 (no RESULTS.md row yet: RESULTS
2026-07-08 logs the full-$\nu$ nullity 1 only). So \textbf{every bit of \Ye{}
evolution comes through the weak columns}, which are 46 of 607 in \msn{80} and
173 of 1518 in \msn{151} \resultsref{2026-07-09}.

Three consequences that shape everything:

\begin{enumerate}
\item \textbf{Invariant \#2 is structural, not stylistic.} Masking a weak
  column as ``equilibrated'' does not lose a little accuracy; it deletes part
  of the only channel that can produce the signal. And weak reactions are
  \emph{not} in detailed balance in this regime --- the inverse of an electron
  capture is neutrino absorption on the daughter, and the neutrinos have
  already left the star (no final-state neutrino blocking;
  \citealt{langanke2000} set $S_\nu = 0$ for exactly this reason); the
  $Z$-restoring channels --- $\beta^{-}$ decay and positron capture --- are
  separate reactions with their own phase space, not the time-reverse of EC.
  There is nothing to balance against. (An earlier version called positron
  capture ``the inverse'' of EC; corrected in the 2026-08-18 audit.)
\item \textbf{$\kap \equiv 1$ on weak columns, structurally}, because the
  engine stores them unpaired ($\fm \equiv 0$). Verified exactly on 8.8
  million / 33.2 million (weak column $\times$ relaxed row) samples with
  $\fp > 0$ \resultsref{2026-07-12}. So weak columns are automatically in any
  $\{\kap > \text{threshold}\}$ active set --- the 95\% \Ye-coverage gate
  passes at 1.0000 \emph{structurally} (\S\ref{sec:gate1}).
\item \textbf{The signal is small and rides on top of large strong-sector
  activity.} The strong sector moves individual $Y_i$ by $O(1)$ relative
  amounts while leaving \Ye{} exactly invariant; the weak sector moves \Ye{} by
  $\sim 10^{-3}$ over a label-length trajectory segment (and by several
  $\times\,10^{-2}$ over the whole Si-burning-to-collapse evolution;
  \citealt{heger2001}). A model that gets the strong sector slightly wrong in a
  way that \emph{leaks} into the charge row destroys the signal. That is
  Tier~2 \S 0.4's entire argument for conservation-by-construction, and this
  is the physical reason it bites.
\end{enumerate}

% =====================================================================
\subsection{Electron capture in a degenerate plasma --- why it is (nearly) a one-way ratchet}
\label{sec:ecratchet}

The rate of electron capture on a nucleus depends on the number of electrons
above the capture threshold. In a degenerate plasma the electrons fill states
up to \eF, so the rate is controlled by whether \eF{} exceeds the threshold
energy $Q_{\mathrm{EC}}$ --- and \eF{} rises with density as $\rho^{1/3}$
(\S\ref{sec:degeneracy}; \citealt[\S IV]{langanke2003a}).
%
\begin{equation}
\lambda_{\mathrm{EC}} \;\propto\; \int_{Q_{\mathrm{EC}}}^{\infty}
\!\!(E - Q_{\mathrm{EC}})^{2}E^{2}\,f(E,\mu_e,T)\,\dd E
\quad\xrightarrow[\text{degenerate}]{}\quad
\lambda \sim \frac{\ln 2}{K}\,\frac{(\eF - Q)^{5}}{(m_ec^2)^5}\ \text{-ish},
\label{eq:lambdaec}
\end{equation}
%
i.e. a steep (roughly fifth-power) sensitivity to how far the Fermi energy
exceeds threshold \derivedhere[{from the LMP phase-space integral ---
\citealt[eq.~5a]{langanke2000}: the exact integrand is
$w\,p\,(Q + w)^{2}\,F(Z, w)\,S_e$; the $\sim E^{5}$ sensitivity is stated in
\citealt[\S IV]{langanke2003a}; it is $(\eF - Q)^{5}$ far above threshold and
$\propto (\eF - Q)^{3}$ just above it; Tier~0 \S 0.1.5 does this properly}].
The practical consequences:

\begin{itemize}
\item \textbf{Rising density switches EC on hard.} As the core contracts, \eF{}
  rises, crosses thresholds one nucleus at a time, and the capture rates jump
  by orders of magnitude (\citealt{suzuki2016}: 3--4 orders over
  $\Delta\log_{10}\rho\Ye = 0.02$ at a threshold). \Ye{} falls faster as the
  collapse proceeds --- the runaway of \S\ref{sec:runaways}.
\item \textbf{The restoring channels --- $\beta^{-}$ decay and positron capture
  --- are suppressed}: $\beta^{-}$ decay is \textbf{Pauli-blocked} in a
  degenerate plasma, since the emitted electron must find an empty state above
  \eF, and there are few \citep{langanke2003a}. So the restoring force is
  suppressed exactly where the driving force is strongest.
\item \textbf{Therefore \Ye{} ratchets downward} --- monotone-non-increasing to
  a good approximation once $\eF \gg$ typical $Q$. During Si burning at
  $\rho \sim 10^{7}$--$10^{9}$~\gcc, however, $\beta^{-}$ decay can balance EC
  for a period (\citealt{heger2001}: $\beta$-decay balances electron capture
  around $\log(t_b - t) = 3.5$--4.5 in the 15~\Msun{} model, and accounts for
  about half of the $\dYe = 0.005$--0.015 by which the new models' central
  \Ye{} at collapse exceeds WW95's), so the strict ratchet is a collapse-phase
  statement; what the gate arithmetic needs is the \emph{sign} of the bias, not
  irreversibility (an earlier version said \Ye{} ``essentially never goes back
  up in this regime''; corrected in the 2026-08-18 audit). Any error in the
  weak sector accumulates in a preferred direction --- which is precisely the
  ``systematic'' row of the accumulation trichotomy (Tier~0 \S IV.3) and the
  reason the per-step gate is set by $N\cdot\delta$ rather than
  $\sqrt{N}\cdot\delta$ (Tier~2 \S 0.3.4).
\end{itemize}

\textbf{And this is where the two networks differ most.} \msn{80} contains only
6 of 9 electron-capture controllers and \textbf{0 of 8} $\beta$-decay partners
of the \Ye-controller set \resultsref{2026-07-08}. It lacks the $\beta^{-}$
partners of the named \Ye{} controllers (it does carry 19 $\beta^{-}$ channels
elsewhere --- RESULTS 2026-07-08 weak-sector breakdown), a more extreme form of
the asymmetry than the physics itself imposes. Tier~0 \S I.4 asks you to
predict the bias direction from that; the answer is that \msn{80} should
over-drive \Ye{} downward relative to \msn{151} in any regime where the
missing $\beta$-decay partners are populated.

That is a \emph{testable prediction about the dataset} which, as far as these
notes can tell, has never been run. Registered here as \reg{T-02}.

% =====================================================================
\subsection{The measured \texorpdfstring{\Ye}{Ye} carriers}
\label{sec:yecarriers}

On the relaxed manifold, in the high-\Ye{} QSE window, the $|\dd\dot Y_e|$
budget is extremely concentrated \resultsref{2026-07-12}:

\begin{center}
\begingroup\small
\begin{tabularx}{\textwidth}{@{}L{1.3cm} L{4.0cm} L{3.0cm} Y L{2.4cm}@{}}
\toprule
\textbf{rank} & \textbf{\msn{80}} & \textbf{share} & \textbf{\msn{151}} &
  \textbf{share} \\
\midrule
1 & \nuc{Ni}{56} EC & 0.62 & \nuc{Ni}{56} EC & 0.28 \\
2--5 & \nuc{S}{31} EC, \nuc{Fe}{52} EC, p EC & $\to$ 0.96 cumulative &
  (ditto class) & $\to$ 0.75 \\
top-20 & & \textbf{1.00} & & \textbf{1.00} \\
\bottomrule
\end{tabularx}
\endgroup
\end{center}

Twenty channels out of 46/173 weak columns carry essentially the entire \Ye{}
signal, and one channel carries a majority of it in \msn{80}. Two design
consequences follow directly and are already in the project's plan:

\begin{itemize}
\item \textbf{Loss weighting} should be concentrated on those channels (a flat
  per-column loss would spend almost all its capacity on columns that do not
  matter for the headline quantity).
\item \textbf{The \nuc{Ni}{56} EC rate's uncertainty is a first-order term in
  the emulator's error budget}, on par with everything the model itself does.
  That is the rate-uncertainty $\to$ \Ye{} propagation flagged as
  \tierone{VIII.C.1} and still open.
\end{itemize}

Note the contrast with the \emph{training} distribution, where the \#1 carrier
is Ne18 $\to$ F18 EC --- a proton-rich fast capture that exists because the
Sobol compositions are random and proton-loaded, not because it happens in a
star \resultsref{2026-07-10}. The \Ye{} physics on the training grid is not the
\Ye{} physics of silicon burning. \S\ref{sec:twodistributions} again.

% =====================================================================
\subsection{Self-check}
\label{sec:selfcheck-weak}

Answer without scrolling up.

\begin{selfcheck}
% source 0.10 Q9
\item Why is $\beta^{-}$ decay unable to undo electron capture in a
  \emph{strongly} degenerate collapsing core --- and when, during Si burning,
  does it manage to? What does that imply about how an emulator's \Ye{} error
  accumulates?
\end{selfcheck}
